\subsection{Définition du projet}
\label{subsec:def_projet}

SP-02 constitue l'un des nombreux projet de fusée développé par l’association AéroIPSA. Ce projet ambitieux vise à concevoir et
réaliser une fusée expérimentale bi-étage, destinée à participer à la campagne C'Space 2026. Il a pour finalité la conception,
la réalisation et la qualification d’un lanceur à deux étages destiné à valider plusieurs technologies innovantes appliquées
aux systèmes de séparation et d’avionique embarquée.

L’objectif principal est de démontrer la faisabilité d’une séparation inter-étage par freinage aérodynamique, alternative aux
méthodes pyrotechniques ou mécaniques conventionnelles. Cette architecture a pour but de simplifier la chaîne d’activation tout
en réduisant les risques liés aux opérations pyrotechniques. Parallèlement, la mission embarque un ensemble de mesures
avioniques dédiées à la caractérisation du vol — accélérations, attitude, pression et altitude — dont les données sont
enregistrées et transmises en temps réel par télémesure vers la station sol.

Le développement du SP-02 s’appuie sur une démarche d’ingénierie système complète, intégrant les volets mécanique, électronique,
logiciel et essais, conformément aux spécifications du cahier des charges \acrshort{CNES}-Planète Sciences pour les fusées
bi-étages. Ce projet vise ainsi à approfondir la maîtrise des processus de conception et de validation expérimentale propres à
l’aéronautique et au spatial, tout en contribuant à l’enrichissement du savoir-faire technique au sein de l’association AéroIPSA.
